\section{Evolución y capacidades centrales de los IDE de IA}

\subsection{Breve historia del IDE}

El IDE (Integrated Development Environment) es el entorno de trabajo central del desarrollo de software: integra en uno solo el editor, el compilador, el depurador y otras herramientas. Su evolución refleja una mejora continua de la eficiencia del desarrollo:

\begin{itemize}
    \item \textbf{1983}: Turbo Pascal, el primer IDE verdadero
    \item \textbf{1997}: Microsoft Visual Studio, con capacidades avanzadas de depuración (C++/Visual Basic)
    \item \textbf{2001}: IntelliJ IDEA, con navegación de código contextual, refactorización y autocompletado de código
    \item \textbf{2015}: VSCode, un editor ligero con un ecosistema de complementos altamente extensible
    \item \textbf{2023}: Cursor, uno de los primeros IDE nativos de IA de uso extendido
\end{itemize}

\begin{knowledgebox}{El «efecto péndulo» en la evolución del IDE}
A lo largo de la evolución del IDE ha existido un «efecto péndulo» (see-saw) entre la integración de funciones y la personalización por parte del desarrollador. La aparición de los IDE de IA vuelve a empujar el péndulo hacia la integración de funciones, incrustando las capacidades de IA profundamente en el entorno de desarrollo.
\end{knowledgebox}

\subsection{Modos de uso de los IDE de IA}

\subsubsection{Modo básico (Bread-and-Butter)}

\begin{itemize}
    \item \textbf{Inline}: autocompletado de código en línea
    \item \textbf{Function}: generación a nivel de función
    \item \textbf{Single-file}: edición de un solo archivo
    \item \textbf{Multi-file}: modificaciones entre archivos
\end{itemize}

\subsubsection{Modo nativo de IA}

\begin{itemize}
    \item \textbf{Background Agents}: ejecución autónoma de tareas en segundo plano
    \item \textbf{Integración con MCP}: conexión con herramientas externas mediante el protocolo MCP
    \item \textbf{Memories}: aprende y recuerda las preferencias del usuario y el contexto del proyecto
    \item \textbf{Bugbot / PR Review}: revisión de código automatizada
\end{itemize}

\subsection{Principios subyacentes de los IDE de IA}

\subsubsection{Tab-complete (autocompletado de código)}

\begin{enumerate}
    \item Se envía cifrada al servidor una ventana de contexto pequeña alrededor del código actual
    \item El servidor ejecuta un LLM de infilling
    \item La sugerencia se devuelve y se muestra en el editor
\end{enumerate}

\subsubsection{Chat (modo de conversación)}

\begin{enumerate}
    \item Los bloques de código se almacenan como embeddings, con lo que se construye un semantic index (usando nombres de archivo ofuscados más el código)
    \item Al consultar el usuario, se recuperan los bloques de código más relevantes y se inyectan como contexto al LLM
    \item El IDE reindexa periódicamente los bloques de código y sincroniza los embeddings
    \item Mediante un Merkle tree se calcula el diff de los bloques de código, lo que permite actualizaciones eficientes
\end{enumerate}

\begin{importantbox}{Aplicación del Merkle tree en la indexación de código}
El Merkle tree es una estructura de árbol de hashes que se usa comúnmente en Git y en blockchain. En los IDE de IA, el Merkle tree se usa para detectar de manera eficiente los cambios en el código: solo los bloques de código que cambiaron necesitan que se recalcule su embedding, con lo que se evita el costo de reconstruir el índice por completo.
\end{importantbox}

\subsection{Resumen de la sección}

El IDE de IA incrusta el LLM profundamente en el flujo de trabajo de desarrollo mediante dos modos: tab-complete y chat. Su funcionamiento subyacente depende de la indexación semántica del código y de un mecanismo eficiente de actualización incremental. Entender estos principios ayuda a aprovechar mejor las capacidades del IDE de IA.
